% GENERATED FILE. Edit canonical lessons, then run npm run generate:books.
% canonical-source-hash: fnv1a64:7bda6ddcb6939f5c
% canonical-lessons: BN-C375-badhyotamulok, BN-C375-onumoti, BN-C375-betikrom, BN-C375-chapano, BN-C375-shithil-kora

\chapter{Must and May}
\label{ch:bn-must-and-may}
\hlchaptermodality{375}

\begin{chapteropening}
\textbf{By the end of this chapter:} \emph{I can say what is required, what is allowed and when a rule bends.}

Complete the last lesson of chapter 375: I can say what is required, what is allowed and when a rule bends.
\end{chapteropening}

\section[\texorpdfstring{\bn{বাধ্যতামূলক} (bādhyôtāmūlôk)}{বাধ্যতামূলক (bādhyôtāmūlôk)}]{\bn{বাধ্যতামূলক} (bādhyôtāmūlôk) — compulsory}

\label{lesson:BN-C375-badhyotamulok}



\begin{quote}
Before the new one: say the Bengali for to control, then the Bengali for a close eye.
\end{quote}

\subsection*{You'll want to know: \bn{বাধ্যতামূলক}}
\textbf{\bn{বাধ্যতামূলক}} — \emph{bādhyôtāmūlôk} — \textquotedblleft{}compulsory\textquotedblright{}.

\begin{quote}
\bn{গাড়ি} \bn{চালালে} \bn{লাইসেন্স} \bn{রাখা} \bn{বাধ্যতামূলক।}
\end{quote}

\begin{quote}
\emph{gāṛi chālāle lāisens rākhā bādhyôtāmūlôk.}
\end{quote}

\begin{quote}
If you drive, you must carry a licence.
\end{quote}

\begin{cousinweb}
\textbf{\bn{বাধ্যতামূলক}} (\emph{bādhyôtāmūlôk}) is a modern Bengali formation: Sanskrit \emph{bādhya}, 'to be pressed or constrained', made abstract with \emph{-tā}, and \emph{mūlaka}, 'rooted in' (Monier-Williams): rooted in obligation, required.
\end{cousinweb}

\begin{grammarlens}[title={What you've built}]
One more word for saying what happens if.
\end{grammarlens}

\subsection*{Your turn}
\begin{itemize}
\raggedright
  \item \emph{Say it:} \emph{bādhyôtāmūlôk}
  \item \emph{Say it:} \emph{bādhyôtāmūlôk}, once more
  \item \emph{Say it:} \emph{gāṛi chālāle lāisens rākhā bādhyôtāmūlôk.} — If you drive, you must carry a licence.
  \item \emph{Recall:} say the Bengali for to control, then the Bengali for a close eye, then say \emph{bādhyôtāmūlôk} again
\end{itemize}

\begin{tcolorbox}[breakable,colback=teal!4,colframe=teal!35!black,title={Before you move on}]
What does \bn{বাধ্যতামূলক} mean? (\textquotedblleft{}Compulsory\textquotedblright{}.) Say it once more.
\end{tcolorbox}
\section[\texorpdfstring{\bn{অনুমতি} (ônumôti)}{অনুমতি (ônumôti)}]{\bn{অনুমতি} (ônumôti) — permission}

\label{lesson:BN-C375-onumoti}



\begin{quote}
Before the new one: say the Bengali for a close eye, then the Bengali for compulsory.
\end{quote}

\subsection*{You'll want to know: \bn{অনুমতি}}
\textbf{\bn{অনুমতি}} — \emph{ônumôti} — \textquotedblleft{}permission\textquotedblright{}.

\begin{quote}
\bn{অনুমতি} \bn{না} \bn{পেলে} \bn{আমরা} \bn{ছবি} \bn{তুলব} \bn{না।}
\end{quote}

\begin{quote}
\emph{ônumôti nā pele āmrā chhôbi tulbô nā.}
\end{quote}

\begin{quote}
We won't take photos unless we get permission.
\end{quote}

\begin{cousinweb}
\textbf{\bn{অনুমতি}} (\emph{ônumôti}) is Sanskrit \emph{anumati}, assent or permission (Monier-Williams). It is asked for with \bn{চাওয়া} and given with \bn{দেওয়া}.
\end{cousinweb}

\begin{grammarlens}[title={What you've built}]
One more word for saying what happens if.
\end{grammarlens}

\subsection*{Your turn}
\begin{itemize}
\raggedright
  \item \emph{Say it:} \emph{ônumôti}
  \item \emph{Say it:} \emph{ônumôti}, once more
  \item \emph{Say it:} \emph{ônumôti nā pele āmrā chhôbi tulbô nā.} — We won't take photos unless we get permission.
  \item \emph{Recall:} say the Bengali for a close eye, then the Bengali for compulsory, then say \emph{ônumôti} again
\end{itemize}

\begin{tcolorbox}[breakable,colback=teal!4,colframe=teal!35!black,title={Before you move on}]
What does \bn{অনুমতি} mean? (\textquotedblleft{}Permission\textquotedblright{}.) Say it once more.
\end{tcolorbox}
\section[\texorpdfstring{\bn{ব্যতিক্রম} (bætikrôm)}{ব্যতিক্রম (bætikrôm)}]{\bn{ব্যতিক্রম} (bætikrôm) — an exception}

\label{lesson:BN-C375-betikrom}



\begin{quote}
Before the new one: say the Bengali for compulsory, then the Bengali for permission.
\end{quote}

\subsection*{You'll want to know: \bn{ব্যতিক্রম}}
\textbf{\bn{ব্যতিক্রম}} — \emph{bætikrôm} — \textquotedblleft{}an exception\textquotedblright{}.

\begin{quote}
\bn{জরুরি} \bn{কাজ} \bn{থাকলে} \bn{ব্যতিক্রম} \bn{হতে} \bn{পারে।}
\end{quote}

\begin{quote}
\emph{jôruri kāj thākle bætikrôm hôte pāre.}
\end{quote}

\begin{quote}
Exceptions can be made if there is urgent work.
\end{quote}

\begin{cousinweb}
\textbf{\bn{ব্যতিক্রম}} (\emph{bætikrôm}) is Sanskrit \emph{vyatikrama}, an overstepping or a reversal of order (Monier-Williams). Bengali uses it for a case that steps outside the rule, an exception.
\end{cousinweb}

\begin{grammarlens}[title={What you've built}]
One more word for saying what happens if.
\end{grammarlens}

\subsection*{Your turn}
\begin{itemize}
\raggedright
  \item \emph{Say it:} \emph{bætikrôm}
  \item \emph{Say it:} \emph{bætikrôm}, once more
  \item \emph{Say it:} \emph{jôruri kāj thākle bætikrôm hôte pāre.} — Exceptions can be made if there is urgent work.
  \item \emph{Recall:} say the Bengali for compulsory, then the Bengali for permission, then say \emph{bætikrôm} again
\end{itemize}

\begin{tcolorbox}[breakable,colback=teal!4,colframe=teal!35!black,title={Before you move on}]
What does \bn{ব্যতিক্রম} mean? (\textquotedblleft{}An exception\textquotedblright{}.) Say it once more.
\end{tcolorbox}
\section[\texorpdfstring{\bn{চাপানো} (chāpāno)}{চাপানো (chāpāno)}]{\bn{চাপানো} (chāpāno) — to impose; to load (onto)}

\label{lesson:BN-C375-chapano}



\begin{quote}
Before the new one: say the Bengali for permission, then the Bengali for an exception.
\end{quote}

\subsection*{You'll want to know: \bn{চাপানো}}
\textbf{\bn{চাপানো}} — \emph{chāpāno} — \textquotedblleft{}to impose; to load (onto)\textquotedblright{}.

\begin{quote}
\bn{আমার} \bn{ওপর} \bn{কাজ} \bn{চাপালে} \bn{আমি} \bn{রাগ} \bn{করব।}
\end{quote}

\begin{quote}
\emph{āmār opôr kāj chāpāle āmi rāg kôrbô.}
\end{quote}

\begin{quote}
If you load work onto me, I'll be angry.
\end{quote}

Past (he/she …): \bn{তিনি} \bn{সবার} \bn{ওপর} \bn{জরিমানা} \bn{চাপালেন}, \emph{tini sôbār opôr jôrimānā chāpālen}.

\begin{cousinweb}
\textbf{\bn{চাপানো}} (\emph{chāpāno}) is the causative of \bn{চাপা} 'to press' (CDIAL 4674, a reconstructed capp, lists Bengali cāpā): to press a load, a tax or a rule down on someone. The one who bears it takes -\bn{এর} \bn{ওপর}.
\end{cousinweb}

\begin{grammarlens}[title={What you've built}]
One more word for saying what happens if.
\end{grammarlens}

\subsection*{Your turn}
\begin{itemize}
\raggedright
  \item \emph{Say it:} \emph{chāpāno}
  \item \emph{Say it:} \emph{chāpāno}, once more
  \item \emph{Say it:} \emph{āmār opôr kāj chāpāle āmi rāg kôrbô.} — If you load work onto me, I'll be angry.
  \item \emph{Recall:} say the Bengali for permission, then the Bengali for an exception, then say \emph{chāpāno} again
\end{itemize}

\begin{tcolorbox}[breakable,colback=teal!4,colframe=teal!35!black,title={Before you move on}]
What does \bn{চাপানো} mean? (\textquotedblleft{}To impose; to load (onto)\textquotedblright{}.) Say it once more.
\end{tcolorbox}
\section[\texorpdfstring{\bn{শিথিল} \bn{করা} (shithil kôrā)}{শিথিল করা (shithil kôrā)}]{\bn{শিথিল} \bn{করা} (shithil kôrā) — to relax a rule; to loosen}

\label{lesson:BN-C375-shithil-kora}



\begin{quote}
Before the new one: say the Bengali for an exception, then the Bengali for to impose.
\end{quote}

\subsection*{You'll want to know: \bn{শিথিল} \bn{করা}}
\textbf{\bn{শিথিল} \bn{করা}} — \emph{shithil kôrā} — \textquotedblleft{}to relax a rule; to loosen\textquotedblright{}.

\begin{quote}
\bn{পরিস্থিতি} \bn{ভালো} \bn{হলে} \bn{সরকার} \bn{নিয়ম} \bn{শিথিল} \bn{করবে।}
\end{quote}

\begin{quote}
\emph{pôristhiti bhālo hôle sôrkār niyôm shithil kôrbe.}
\end{quote}

\begin{quote}
If the situation improves, the government will relax the rules.
\end{quote}

Past (he/she …): \bn{তিনি} \bn{নিয়মটা} \bn{শিথিল} \bn{করলেন}, \emph{tini niyômṭā shithil kôrlen}.

\begin{cousinweb}
\textbf{\bn{শিথিল} \bn{করা}} (\emph{shithil kôrā}) is Sanskrit \emph{śithila}, loose or slack (Monier-Williams), with \bn{করা}: to let a rule or a grip go slack.
\end{cousinweb}

\begin{grammarlens}[title={What you've built}]
One more word for saying what happens if.
\end{grammarlens}

\subsection*{Your turn}
\begin{itemize}
\raggedright
  \item \emph{Say it:} \emph{shithil kôrā}
  \item \emph{Say it:} \emph{shithil kôrā}, once more
  \item \emph{Say it:} \emph{pôristhiti bhālo hôle sôrkār niyôm shithil kôrbe.} — If the situation improves, the government will relax the rules.
  \item \emph{Recall:} say the Bengali for an exception, then the Bengali for to impose, then say \emph{shithil kôrā} again
\end{itemize}

\begin{tcolorbox}[breakable,colback=teal!4,colframe=teal!35!black,title={Before you move on}]
What does \bn{শিথিল} \bn{করা} mean? (\textquotedblleft{}To relax a rule; to loosen\textquotedblright{}.) Say it once more.
\end{tcolorbox}
